% Audited mathematical working proof: root, four_regular, and path_stability passed full independent line audits.
\section{A bootstrap to the unrestricted sharp-order error}
\label{sec:unrestricted-sharp-bootstrap}

This audited mathematical working proof combines the excess-free girth-raising theorem with two conditional clean-block estimates.
The original qualitative marked-cycle estimate suffices through localization and affine-forest bounds.
The independently established polynomial forest theorem supplies a stronger version of that input.
Thus the argument does not assume the sharp conclusion that it seeks to prove.
All logarithms have base two.

\begin{theorem}[An exact half bound above a fixed girth]
\label{thm:bootstrap-exact-half-forest}
There are absolute integers $D_*,R$ such that the following holds.
Let $G$ be a connected cubic multigraph of rank $r\ge R$.
Let $P$ be connected or empty, with $|P|\le3\log r$.
If
\[
 \operatorname{girth}(G-P)\ge D_*,
\]
then
\[
 \Pr(\xi_G-P\text{ is a forest})\le\frac12.
\]
Here an acyclic graph has infinite girth.
\end{theorem}

\begin{proof}
Take $D_*$ from Theorem~\ref{thm:four-excess-free-girth-raising}.
If $P$ is empty, the only even forest is the empty state, of probability $2^{-r}$.
Assume $P$ is nonempty and, for contradiction, that the forest probability exceeds one half.
Write $x=\log r$, $z=\log x$, and put
\[
 A=6e^{168},\qquad
 L=\lfloor Ax\rfloor,\qquad
 t=\lfloor z-2\log z\rfloor.
\]
Use the common forward dyadic sequence $2^jD_*$.
Let $D_0$ be its first value at least $\sqrt x$, and let $D_1$ be its first value at least $x/16$.
For sufficiently large $r$,
\[
 \sqrt x\le D_0<2\sqrt x,\qquad
 x/16\le D_1<x/8,\qquad D_*<D_0<D_1.
\]
Apply the maximal-packing recurrence from the girth-raising theorem along this one sequence.
It supplies nested sets $S_0\subseteq S_1$, disjoint from $P$, such that
\[
 J_j=G-(P\cup S_j)
\]
has girth greater than $D_j$, and
\begin{align*}
 \log(|S_0|+1)
 &\le2D_0+K\log^2(D_0+2)=o(x),\\
 \log(|S_1|+1)
 &\le2D_1+K\log^2(D_1+2)
 \le x/4+o(x).
\end{align*}
In particular,
\[
 |S_0|=r^{o(1)},\qquad |S_1|\le r^{1/4+o(1)}=o(r/x).
\]
Each $S_j$ is the vertex union of a maximal disjoint family of cycles of length at most $D_j$ outside $P$:
the construction only adds disjoint cycles, and its final complement contains no further such cycle.

We use the non-backtracking trace and local cycle-weight estimates from Section~\ref{sec:polynomial-foundations}.
The link and moment estimates are Lemmas~\ref{lem:sharp-uniform-links} and~\ref{lem:sharp-half-moment}.
Their required parameters are verified here.
Write
\[
 s_1=\lfloor(D_1-1)/2\rfloor\ge x/32-2.
\]
The graph $J_1$ is simple, has order $v_1=2r+o(r)$, and has degree deficit
$d_1\le3(|P|+|S_1|)$.
Consequently $Ld_1/(2v_1)=o(1)$ and $L^2 2^{-s_1}=o(1)$.
The even trace sum over $6x\le\ell\le L$ therefore gives
\[
 \mu_1:=
 \sum_{\substack{C\subseteq J_1\\|C|\le L}}2^{-|C|}
 \ge\left(\frac14-o(1)\right)\ln(A/6)>20.
\]

The large-scale link estimate applies because $8\le D_1\le L$ and
$2(t+2\lfloor D_1/8\rfloor)\le D_1$ for large $r$.
The marked set is connected, so it gives the bound
\[
 b_1\le
 1+\frac{2L}{t+1}+2^{t+1}+578(L/D_1)^2
 \le\frac{3Ax}{z}
\]
on the number of links between disjoint cycles of length at most $L$ in $J_1$.
The exceptional-pair error is
\[
 \eta_1\le\frac32\,2^{b_1}L^3 2^{-s_1}
 \le r^{-1/32+o(1)}=o(1).
\]
Only the first refined moment conclusion is needed:
\[
 W_{2,1}>\mu_1^2-(1+\eta_1)\mu_1.
\]
It follows from forest probability greater than one half, without a lower bound on that excess.
No moment estimate is used at scale $D_0$.

Apply Lemma~\ref{lem:sharp-nested-chain}, with the connected set $P$ as anchor.
It gives an auxiliary cycle $C\subseteq J_1$ of length at most $L$, a neighbor $u$ of $C$, and a nested chain
$C_1,\ldots,C_T$ of vertex-disjoint cycles in the component of $G-V(C)$ containing $P$.
Each chain cycle separates $u$ from $P$, and
\[
 T\ge\frac{8\,2^{s_1}}{L^3}
 \ge\frac{2}{A^3}\frac{r^{1/32}}{x^3}.
\]
For reference, the moment inequality leaves partner weight at least eight after discarding ordered far pairs and adjacent pairs.
Choosing one of at most $L$ neighbors of $C$ leaves weight at least $8/L$.
The local cycle-weight bound $L2^{-s_1}$ and maximal vertex-disjointness give the displayed chain length.

Let $H=G-V(C)$, and let $U_i$ be the $u$-component of $H-V(C_i)$.
The nesting gives $U_i\subseteq U_{i+1}$ and $C_i\subseteq U_{i+1}$.
Put $q=\lceil128A^2x\rceil+2$.
Choose disjoint index intervals $[a,b]$ with $b-a=q$, leaving a separator between successive intervals.
For each one, take the full block
\[
 R(a,b)=U_b\setminus\bigl(U_a\cup V(C_a)\bigr).
\]
These blocks are pairwise disjoint.
Each contains all $q-1$ intermediate chain cycles.
Since each of those cycles has length greater than $D_1$,
\[
 |R(a,b)|\ge(q-1)(D_1+1)\ge8L^2
\]
for sufficiently large $r$.
Every edge of $H$ leaving a block ends on $C_a$ or $C_b$, so its boundary in $H$ has at most $2L$ edges.

Discard blocks meeting a neighbor of $C$, a vertex of $S_0$, or a vertex of $P$.
This removes at most $L+|S_0|+1$ blocks.
The last bound uses connectedness of $P$: it cannot cross a block boundary, since that boundary ends on chain cycles disjoint from $P$.
The total number of blocks is of order at least $r^{1/32}/x^4$.
Thus at least two remain after the $r^{o(1)}$ discards.

Every retained block avoids $P\cup S_0$, has girth greater than $D_0$, and has boundary at most $2L$ in $G$.
The boundary statement follows because the block meets no neighbor of the removed auxiliary cycle $C$.
Finally,
\[
 \frac{D_0}{\log(L+2)}\longrightarrow\infty.
\]
Proposition~\ref{prop:path-two-clean-blocks} therefore bounds the original forest probability by $9/25$.
This contradicts the assumption that it exceeds one half and proves the theorem.
\end{proof}

\begin{theorem}[The unrestricted sharp-order upper bound]
\label{thm:bootstrap-unrestricted-sharp-error}
There is an absolute constant $K$ such that every finite graph of cycle rank $r\ge1$ satisfies
\[
 c(G)\le2^{r-1}+Kr2^{r/2}.
\]
The order of the second term is necessary on an infinite sequence of ranks.
\end{theorem}

\begin{proof}
Join distinct components by bridges.
This preserves the cycle count and the rank $r=|E|-|V|+\operatorname{comp}$, so it suffices to prove the connected case.

We use two established inputs.
First, the extremal cycle count $\psi(r)$ is attained by a finite connected cubic multigraph for $r\ge2$.
Attainment starts from a finite graph because the nonempty set of possible cycle counts is a set of integers bounded by $2^r-1$.
The reduction deletes leaves, suppresses degree-two vertices, and replaces each higher-degree vertex by a trivalent tree.
It preserves rank and injects the original cycles.
Multigraphs and simple graphs have the same extremal value because subdivision preserves rank and cycle count.
Second, Theorem~\ref{thm:sharp-order-four-cut-error} gives the claimed order, with an absolute constant $K_4$, for loopless cubic multigraphs with a cyclic cut of size at most four.

Choose an absolute rank threshold $R_0$ above that in
Theorem~\ref{thm:bootstrap-exact-half-forest} and above the fixed-girth deficit thresholds for every girth below $D_*$.
The latter thresholds exist by Corollary~\ref{cor:path-unbounded-girth-reduction}.
Increase $R_0$ if needed for the cubic Moore bound $g\le3\log r$.
Choose $K\ge K_4$ large enough that the asserted bound holds for every rank below $R_0$.
This is possible using the elementary bound $\psi(r)\le2^r-1$.
We prove the theorem by strong induction on the rank.

Let $r\ge R_0$, and choose a cubic multigraph attaining $\psi(r)$.
If it has a loop $e$, that loop is the only cycle containing $e$.
Deleting $e$ leaves a connected graph of rank $r-1$, so the induction hypothesis gives
\[
 c(G)=c(G-e)+1
 \le2^{r-2}+K(r-1)2^{(r-1)/2}+1
 \le2^{r-1}+Kr2^{r/2}.
\]
Assume henceforth that $G$ is loopless.

If it has a cyclic cut of size at most four, apply the established bound with constant $K_4$.
Otherwise its girth is at least five at sufficiently large rank:
a cycle of length at most four has boundary at most four and leaves a cyclic exterior once the rank exceeds an absolute constant.
If its girth is below $D_*$, the fixed-girth deficit gives
$c(G)<2^{r-1}$.
These cases are complete.

In the remaining case, choose a shortest cycle $W$.
It has at most $3\log r$ vertices, and $G-V(W)$ has girth at least $D_*$.
Theorem~\ref{thm:bootstrap-exact-half-forest}, with $P=V(W)$, gives forest probability at most one half.
For any edge $e$ of $W$, each cycle containing $e$ gives two distinct even states by toggling $W$.
Both restrictions outside $V(W)$ are forests.
Consequently the number $c_e$ of cycles containing $e$ satisfies
\[
 2c_e\le\frac12\,2^r,
 \qquad c_e\le2^{r-2}.
\]
The graph $G-e$ is connected and has rank $r-1$.
Apply the induction hypothesis to this general connected graph:
\[
 c(G)=c(G-e)+c_e
 \le2^{r-1}+K(r-1)2^{(r-1)/2}
 \le2^{r-1}+Kr2^{r/2}.
\]
This completes the induction.
It does not require $G-e$ to remain cubic.

Finally, the cyclic-square family from the sharp four-cut theorem has ranks $r=2q+1$ and excess
\[
 \left(2q+\frac12\right)2^q-3q
 =\Theta(r2^{r/2}).
\]
Thus the error term cannot be replaced uniformly by $o(r2^{r/2})$.
\end{proof}
